% Generated draft: requires mathematical and source review.
\documentclass[11pt]{article}
\usepackage[margin=25mm]{geometry}
\usepackage{fontspec}
\setmainfont{Latin Modern Roman}
\usepackage{amsmath,amssymb,mathtools,unicode-math}
\setmathfont{Latin Modern Math}
\usepackage{enumitem,needspace,xurl,hyperref}
\setlength{\parindent}{0pt}
\setlength{\parskip}{5pt}
\setlength{\emergencystretch}{3em}
\newcommand{\N}{\mathbb{N}}
\newcommand{\Z}{\mathbb{Z}}
\newcommand{\Q}{\mathbb{Q}}
\newcommand{\R}{\mathbb{R}}
\newcommand{\C}{\mathbb{C}}
\newcommand{\F}{\mathbb{F}}
\newcommand{\A}{\mathbb{A}}
\newcommand{\PP}{\mathbb{P}}
\newcommand{\E}{\mathbb{E}}
\newcommand{\eps}{\varepsilon}
\newcommand{\dd}{\,\mathrm{d}}
\begin{document}
\section*{UnsolvedMath ID: 30001066}
\textbf{Source problem code:} catalogue entry
\subsection{Definitions and mathematical statement}
Let \(\mathcal K=\{K_1,\ldots,K_m\}\) be a finite family of pairwise disjoint convex subsets of \(\R^d\). A line transversal is an affine line meeting every \(K_i\). Give the space of affine unoriented lines its usual topology, and put
\[
\mathcal T(\mathcal K)=\{\ell:\ell\cap K_i\ne\varnothing\text{ for every }i\}.
\]
A transversal \(\ell\) is isolated if there is an open neighborhood \(U\) in this line space with \(U\cap\mathcal T(\mathcal K)=\{\ell\}\). This is local isolation, not necessarily uniqueness among all transversals.

Does there exist a finite bound \(h(d)\), depending only on the ambient dimension, such that for every such finite family and every isolated transversal \(\ell\), some subfamily \(\mathcal K'\subseteq\mathcal K\) with \(|\mathcal K'|\leq h(d)\) already has \(\ell\) as an isolated transversal? In particular, can one take \(h(d)=2d-1\), as suggested by the source's proposed analogue of the theorem for disjoint balls?

No compactness, closedness, full-dimensionality, smoothness, or strict convexity of the arbitrary convex sets is stated. The desired uniform bound and the specific \(2d-1\) proposal are separate levels of the question. No regular-body hypothesis or claimed bound is added to the unrestricted target.
\subsection{Short English statement}
Ask whether local isolation of a line meeting disjoint convex sets is witnessed by a dimension-bounded subcollection, with \(2d-1\) as the proposed bound. Isolation need not mean global uniqueness, and arbitrary convex sets remain the stated target.
\subsection{Sources}
\begin{itemize}
\item Local UnsolvedMath registry, permanent ID \(30001066\), source code \texttt{catalogue entry}, ``Helly Numbers for Isolated Line Transversals.'' Registry curation: Ulam AI / UnsolvedMath Contributors, CC BY 4.0; underlying sources retain their supplied provenance.
\item \emph{Discrete Geometry} (2009), Oberwolfach Reports, pp.\ 2477--2556, DOI \url{https://doi.org/10.4171/owr/2008/44}; supplied literature label ``OWR Report 44.''
\item Supplied ``Helly-type transversal results'': \url{https://www.ens-lyon.fr/LIP/Arenaire/SYMB/teams/vegas/vegas2.pdf}.
\item Supplied literature triage, \texttt{literature\_checked\_at}: 2026-08-21: labeled open; says the report proves the sharp bound for disjoint balls and selected regular classes, while extension to arbitrary disjoint convex sets remains open. This is supplied dated curation, not independently verified; the exact regularity hypotheses and known bound need source review.
\end{itemize}
\end{document}
